% Shared academic preamble. White page, black text, serif body.
\documentclass[a4paper,11pt,fontset=none]{ctexart}
\setmainfont{Noto Serif}
\setCJKmainfont{Noto Serif CJK SC}[BoldFont={Noto Serif CJK SC Bold}]
\setCJKsansfont{Noto Sans CJK SC}[BoldFont={Noto Sans CJK SC Bold}]
\setCJKmonofont{Noto Sans CJK SC}
\usepackage{geometry}
\geometry{a4paper,margin=2.3cm,headheight=14pt}
\usepackage{booktabs}
\usepackage{longtable}
\usepackage{array}
\usepackage{tabularx}
\usepackage{tikz}
\usetikzlibrary{arrows.meta,positioning}
\usepackage{float}
\usepackage{caption}
\captionsetup{font=small,labelfont=bf,skip=6pt}
\usepackage{fancyhdr}
\usepackage{setspace}
\setstretch{1.15}
\usepackage{enumitem}
\setlist[itemize]{itemsep=0.2em,topsep=0.35em,leftmargin=1.6em}
\setlist[enumerate]{itemsep=0.2em,topsep=0.35em,leftmargin=1.6em}
\usepackage{amsmath}
\usepackage{xurl}
\usepackage[hidelinks]{hyperref}
\pagestyle{fancy}
\fancyhf{}
\fancyhead[L]{\small 盼盼法式小面包}
\fancyhead[R]{}
\fancyfoot[C]{\thepage}
\renewcommand{\headrulewidth}{0.4pt}
\setcounter{secnumdepth}{3}
\setcounter{tocdepth}{2}
\emergencystretch=2em
\Urlmuskip=0mu plus 1mu
\tikzset{
  layer/.style={draw, align=flush center, font=\small, inner sep=5pt, text width=3.35cm, minimum height=1.05cm},
  wide/.style={draw, align=flush center, font=\small, inner sep=6pt, text width=14.2cm},
  flow/.style={-{Stealth[length=2.2mm]}, semithick},
  abox/.style={draw, fill=white, align=flush center, font=\scriptsize, inner sep=2.5pt,
    text width=27mm, minimum height=11mm},
  sbox/.style={draw, dashed, fill=white, align=flush center, font=\scriptsize, inner sep=2.5pt,
    text width=24mm, minimum height=11mm},
  alab/.style={font=\scriptsize\bfseries, align=center, text width=12mm},
  aflow/.style={-{Stealth[length=1.6mm]}, semithick},
  aside/.style={-{Stealth[length=1.6mm]}, densely dashed, semithick}
}
\newcommand{\contesttitlepage}[6]{%
  \begin{titlepage}
  \thispagestyle{empty}
  \centering
  \vspace*{1.5cm}
  {\large 首届江苏省“人工智能+”创新应用大赛\par}
  \vspace{0.55cm}
  {\large #1\par}
  \vspace{1.35cm}
  {\rule{0.42\textwidth}{0.5pt}\par}
  \vspace{0.75cm}
  {\Huge\bfseries #2\par}
  \vspace{0.4cm}
  {\Large #3\par}
  \vspace{0.95cm}
  {\LARGE #4\par}
  \vspace{0.5cm}
  {\large #5\par}
  \vspace{0.75cm}
  {\rule{0.42\textwidth}{0.5pt}\par}
  \vfill
  {\large #6\par}
  \vspace*{1.1cm}
  \end{titlepage}
}

\hypersetup{pdftitle={苏算绿枢 佐证材料},pdfauthor={盼盼法式小面包}}
\begin{document}
\contesttitlepage
  {“三苏”创新应用赛道\\赛题 AI-JIANGSU-Q36}
  {苏算绿枢}
  {SuCompute-Opt}
  {佐证材料}
  {仿真评测、权重说明与边车部署方案}
  {团队：盼盼法式小面包\\2026 年 11 月\\复现主机 4 vCPU\\Python 3.12.3，torch 2.7.0+cpu}

\begin{abstract}
评测在 \path{02-sansu-computing-dispatch} 目录执行，使用仓库内已训练权重，没有重训。种子 313，四场景各 168 小时。本次滚动规划最长 0.0957 秒，SuCompute 故障编排最长 0.0629 秒。脚本原图 \path{comparison.png} 与 \path{training.png} 仍在 \path{deliverables/repro/}，本文以表列出同一批读数，不嵌栅格图。小时级 Gym 的加权 PUE 与碳排按 \path{baseline_seed313.md} 分列。
\end{abstract}
\tableofcontents
\newpage

\section{两次命令，两套不能混用的模型}
依赖见 \path{requirements.txt}。三条命令是：
\begin{quote}\small
\path{python weights/verify_weights.py}\\
\path{python scripts/run_baselines.py --seed 313}\\
\path{python benchmarks/run_benchmark.py --seed 313 --workers 2 --out} <输出目录>
\end{quote}
第一条检查五份权重的 SHA-256、ONNX 与 PyTorch logits，并做一次卷积前向。第二条只跑小时级 Gym 的 random 和 least-loaded greedy。第三条才产生 RR-Local、Least-Loaded、Cost-First、SuCompute、MILP-only、Safe-RL-only。

\begin{table}[ht]
\centering
\caption{两套沙盘不能混用。}
\label{tab:two}
\small
\begin{tabularx}{\textwidth}{lXXXX}
\toprule
沙盘 & 步长 & 功率系数 & 碳因子 & 策略 \\
\midrule
\path{js_computesim} & 1 小时 & 1{,}000 kW/EFlops & 分省年平均 & random、greedy \\
backend 五分钟引擎 & 5 分钟 & 10{,}000 kW/EF & 演示值 0.55 & 六策略，含三条基线 \\
\bottomrule
\end{tabularx}
\end{table}

本次 verify 输出：Verified 5 hashes, 38987 bytes; max logit error 3.81e-06; forecast (24, 16, 2)。运行环境见 \path{results_summary.json}：Python 3.12.3，x86\_64，scipy 1.15.2，torch 2.7.0+cpu。

逐小时 CSV 的 SHA-256 本次为 \path{034b8c871e44310fd2f1c5131979a8ad1418e328e1577e137b6c8e3b486626d7}。仓库冻结 \path{benchmarks/results_detail.csv} 为 \path{07625d6e37ac223cf1d6b936b6cc3b972d0b16c0ceb32d21cefc413e09437f7d}。周报表四舍五入一致，但哈希不同，不以冻结哈希冒充本次复现。

\section{六策略周指标}
四个场景等权。有效在线利用率等于按 \(\eta\) 折算的有效工作量除以含保活的已通电容量时。SLA 分母是到达任务数，拒单计入违约。国产有效份额的约束下限是 30\%，表中是实际份额。高峰时段为每日 09--11 与 17--19 时到达的任务。碳强度单位是 kgCO\textsubscript{2}/EF$\cdot$h。

\begin{table}[ht]
\centering
\caption{四场景等权均值。读数与 deliverables/repro/results\_report.md 一致。}
\label{tab:six}
\footnotesize
\setlength{\tabcolsep}{2.6pt}
\begin{tabular}{lrrrrrrrr}
\toprule
策略 & 利用率 & SLA & kWh/EF$\cdot$h & PUE & 绿电 & 碳强度 & 高峰 SLA & 国产份额 \\
\midrule
RR-Local & 62.72\% & 95.72\% & 17705.0 & 1.3760 & 25.26\% & 7277.5 & 92.64\% & 56.30\% \\
Least-Loaded & 60.92\% & 98.13\% & 18181.2 & 1.3497 & 25.98\% & 7401.3 & 97.04\% & 54.79\% \\
Cost-First & 68.98\% & 94.97\% & 16271.5 & 1.3093 & 27.26\% & 6509.9 & 93.69\% & 56.93\% \\
SuCompute & 80.01\% & 99.83\% & 13945.7 & 1.2957 & 35.95\% & 4911.9 & 99.80\% & 52.19\% \\
MILP-only & 74.10\% & 99.83\% & 14298.4 & 1.3177 & 35.57\% & 5066.2 & 99.81\% & 52.20\% \\
Safe-RL-only & 80.14\% & 99.81\% & 13954.4 & 1.2965 & 35.71\% & 4933.7 & 99.76\% & 51.83\% \\
\bottomrule
\end{tabular}
\end{table}

\path{benchmarks/run_benchmark.py} 另写 \path{comparison.png}。六组柱从左到右为 RR-Local、Least-Loaded、Cost-First、SuCompute、MILP-only、Safe-RL-only。横轴四个场景为常规工作日、绿电大发、尖峰保电、故障压测。图中数字与表~\ref{tab:six} 同源，本文不嵌该栅格图。

\section{相对差和脚本里的下限}
利用率、能耗、绿电、碳强度是相对变化。SLA 是百分点差。PUE 是绝对差。能耗列的正号表示苏算绿枢更低。

\begin{table}[ht]
\centering
\caption{SuCompute 相对各对照的变化。}
\label{tab:delta}
\footnotesize
\setlength{\tabcolsep}{4pt}
\begin{tabular}{lrrrrrr}
\toprule
对照 & 利用率 & SLA 百分点 & 能耗 & PUE & 绿电 & 碳强度 \\
\midrule
RR-Local & +27.57\% & +4.12 & +21.23\% & +0.0803 & +42.31\% & +32.51\% \\
Least-Loaded & +31.33\% & +1.70 & +23.30\% & +0.0540 & +38.36\% & +33.63\% \\
Cost-First & +15.99\% & +4.87 & +14.29\% & +0.0136 & +31.90\% & +24.55\% \\
MILP-only & +7.97\% & +0.01 & +2.47\% & +0.0220 & +1.08\% & +3.05\% \\
Safe-RL-only & $-0.16\%$ & +0.03 & +0.06\% & +0.0007 & +0.67\% & +0.44\% \\
\bottomrule
\end{tabular}
\end{table}

\begin{table}[ht]
\centering
\caption{“达到”是脚本相对 RR-Local 的下限判断，作用域是这次合成周。最后一行是本机墙钟，不是冻结报告里的 0.0891 秒与 0.0803 秒。}
\label{tab:gate}
\begin{tabular}{lrl}
\toprule
脚本中的下限 & 本次印出值 & 脚本结论 \\
\midrule
利用率 +15\%～22\% & 27.57\% & 达到 \\
各场景 SLA \(\ge\) 99.2\% & 99.80\% & 达到 \\
单位能耗降低 12\%～18\% & 21.23\% & 达到 \\
有效在线利用率 \(\ge\) 78\% & 80.01\% & 达到 \\
绿电相对提升 18\%～25\% & 42.31\% & 达到 \\
碳强度降低 15\%～22\% & 32.51\% & 达到 \\
加权 PUE 降低 \(\ge\) 0.06 & 0.0803 & 达到 \\
最长故障编排 \(\le\) 3 秒 & 0.0629 秒 & 达到 \\
\bottomrule
\end{tabular}
\end{table}

\section{移峰、重排和本机墙钟}
MILP-only 保留同一套滚动计划和重排，用配额贪心代替 PPO。Safe-RL-only 加载同一个 PPO，但没有预报、移峰和逐小时重排。批任务“推迟至少 2 小时”只统计预训练和微调里首次开工推迟不少于 24 个五分钟槽的作业。

\begin{table}[ht]
\centering
\caption{Safe-RL-only 的推迟和重排为 0，因为它不移动批任务的开工时刻。全装机口径的分母含未通电容量，所以远低于有效在线利用率。}
\label{tab:shift}
\footnotesize
\setlength{\tabcolsep}{3pt}
\begin{tabular}{llrrrr}
\toprule
场景 & 策略 & 推迟 \(\ge\) 2h & 重排次数 & 高峰 SLA & 全装机有效利用率 \\
\midrule
normal\_weekday & SuCompute & 3305 & 1354 & 99.83\% & 32.38\% \\
normal\_weekday & MILP-only & 3060 & 2334 & 99.83\% & 32.33\% \\
normal\_weekday & Safe-RL-only & 0 & 0 & 99.78\% & 32.36\% \\
high\_renewable & SuCompute & 3485 & 1152 & 99.76\% & 32.34\% \\
high\_renewable & MILP-only & 3144 & 1967 & 99.78\% & 32.33\% \\
high\_renewable & Safe-RL-only & 0 & 0 & 99.73\% & 32.34\% \\
peak\_shaving & SuCompute & 3508 & 1470 & 99.82\% & 34.27\% \\
peak\_shaving & MILP-only & 3327 & 2467 & 99.84\% & 34.22\% \\
peak\_shaving & Safe-RL-only & 0 & 0 & 99.78\% & 34.20\% \\
fault\_stress & SuCompute & 3299 & 1402 & 99.80\% & 32.34\% \\
fault\_stress & MILP-only & 3051 & 2475 & 99.80\% & 32.31\% \\
fault\_stress & Safe-RL-only & 0 & 0 & 99.74\% & 32.35\% \\
\bottomrule
\end{tabular}
\end{table}

本次全体滚动规划最长 0.0957 秒，出现在故障场景的 MILP-only。SuCompute 在故障场景的编排最长 0.0629 秒，\path{solve_max_s} 为 0.0681 秒，\path{measured_solves = 168}。故障轨迹的 \path{elapsed_s} 含计划重算和内存中的切流，不含生产环境的检查点字节复制。

仓库冻结 \path{results_report.md} 写的是规划 0.0891 秒、故障编排 0.0803 秒。那是生成冻结包的机器的墙钟。本机是 4 vCPU 的 Intel Xeon 虚拟机，两个数都变了，两个都仍小于脚本里的 3 秒和 1 秒门槛。引用时间以本页为准。

\section{小时级 Gym 是另一张表}
\path{scripts/run_baselines.py} 不产生 RR-Local、Cost-First 或 SuCompute。random 与 greedy 都不看电价和绿电，所以绿电大发日的放置与常规工作日相同，只有碳和绿电份额变化。绝对 MWh、吨和万元使用 1{,}000 kW/EFlops，脚本原文写明这不是全省能源平衡。

表~\ref{tab:gym} 按 \path{deliverables/repro/baseline_seed313.md} 分列 \path{weighted_pue} 与 \path{carbon_t}。先前排版稿把碳排吨数印在“加权 PUE”列下，此处改正。完成比、SLA 违约、平均利用率和绿电份额的数字没有改。

\begin{table}[ht]
\centering
\caption{完成比是 served\_ratio。SLA 违约是 sla\_violation\_ratio，方向与五分钟引擎的 SLA 达标率相反。碳排单位为吨。}
\label{tab:gym}
\footnotesize
\setlength{\tabcolsep}{3pt}
\begin{tabular}{llrrrrrr}
\toprule
场景 & 策略 & 完成比 & SLA 违约 & 平均利用率 & 加权 PUE & 碳排（t） & 绿电份额 \\
\midrule
weekday & random & 0.8109 & 0.1891 & 0.8873 & 1.3918 & 37615.8845 & 0.0280 \\
weekday & greedy & 0.9404 & 0.0596 & 0.8464 & 1.3903 & 36517.3900 & 0.0288 \\
green\_surplus & random & 0.8109 & 0.1891 & 0.8873 & 1.3918 & 36144.9595 & 0.0660 \\
green\_surplus & greedy & 0.9404 & 0.0596 & 0.8464 & 1.3903 & 35046.4324 & 0.0679 \\
peak\_conservation & random & 0.6874 & 0.3126 & 0.9065 & 1.3927 & 36021.1673 & 0.0292 \\
peak\_conservation & greedy & 0.8043 & 0.1957 & 0.8921 & 1.3925 & 35810.8230 & 0.0294 \\
node\_failure & random & 0.8049 & 0.1951 & 0.8913 & 1.4077 & 37415.1237 & 0.0281 \\
node\_failure & greedy & 0.9337 & 0.0663 & 0.8541 & 1.4067 & 36451.9217 & 0.0289 \\
\bottomrule
\end{tabular}
\end{table}

常规工作日 greedy 的完成比是 0.9404，违约 0.0596；random 的完成比是 0.8109，违约 0.1891。尖峰保电日 greedy 违约升到 0.1957。这些数只说明该 Gym 在种子 313 下的放置，不能拿来和支持卷积预报的五分钟 SLA 99.83\% 相减。

小时环境的预测接口当前等于生成器真值。把它换成真实预报器时，替换 \path{planning_snapshot()} 的到达数组即可。这次基线复现没有做这件替换。同一文件中的 \path{energy_mwh}、\path{cost_10k_cny}、\path{cross_domain_eflops} 与 \path{reward} 未列入上表，仍以该 Markdown 为准。

\section{权重包是时间卷积和小型 PPO}
\path{manifest.json} 的 version 为 2。预报器是季节残差 TCN：Conv1d(4\(\to\)24, \(k=5\))、Conv1d(24\(\to\)16, \(k=3\), dilation\(=2\))、Conv1d(16\(\to\)2, \(k=1\))，残差接在滞后 24 小时与 48 小时的平均上，输出经 clamp 不低于 0。策略网是两层 Tanh 的候选打分器和价值头，输入 15 维特征，部署前还要加上 \path{dispatch_bias} 和可行性掩码。

\begin{table}[ht]
\centering
\caption{五份文件合计 38987 字节，低于仓库设定的 15 MB 上限。SHA-256 前 16 位与 manifest.json 一致。}
\label{tab:w}
\small
\begin{tabularx}{\textwidth}{lrlX}
\toprule
文件 & 字节 & SHA-256 前 16 位 & 说明 \\
\midrule
policy.pt & 9{,}669 & fc11e917392a9140 & PPO actor and critic state\_dict \\
forecast.pt & 9{,}621 & 4153bc439493ff61 & seasonal residual TCN state\_dict \\
policy.onnx & 5{,}437 & 2f7b6d8032225469 & portable actor logits, opset 17 \\
forecast\_metrics.json & 707 & 5ccb04478642254c & independent train/validation MAE and positive-target MAPE \\
ppo\_training.csv & 8{,}539 & 6d3543ee4dd1219e & on-policy sample reward, violation rate and projected multiplier \\
\bottomrule
\end{tabularx}
\end{table}

本次校验：Verified 5 hashes, 38987 bytes; max logit error 3.81e-06; forecast (24, 16, 2)。ONNX 只导出策略 logits，预报仍用 PyTorch。

特征顺序与代码 \path{FEATURES} 一致：eta、tariff、green\_fraction、pue\_excess、pool\_fill、marginal\_\-powered\_\-efh\_\-per\_\-efh、latency\_ratio、cross、domestic、quota、sla\_risk、wait\_\-fraction\_\-of\_\-6h、energy\_per\_work、available、urgent。

不要把这套权重描述成 PatchTST 或 LightGBM。\path{research.md} 和 \path{solution.md} 里出现过这两个名字，是选型笔记，不是已发布权重。重新训练的入口是 \path{benchmarks/train_forecast.py} 与 \path{train_policy.py}，训练后再跑 \path{verify_weights.py --release}。本次没有执行训练。

\section{预报误差}
验证种子 72。评测周使用种子 313，不在这张误差表里。\path{run_benchmark.py} 按 \path{weights/ppo_training.csv} 绘制训练图：绿线是 10 次更新的滑动平均，红虚线是违约率 0.5\%。该图不嵌入本文。

\begin{table}[ht]
\centering
\caption{训练集与验证集。验证集到达 MAPE 43.96\% 说明卷积只抓住了日周期的一部分。周评测里的移峰仍然使用这份预报，因此利用率优势不能解释成预报已经准确。}
\label{tab:mape}
\begin{tabular}{llrrr}
\toprule
划分 & 变量 & MAE & MAPE & 非零样本 \\
\midrule
训练 & 到达（作业/小时） & 2.042 & 43.71\% & 137{,}010 \\
训练 & 绿电比例 & 0.0198 & 12.43\% & 152{,}064 \\
验证 & 到达（作业/小时） & 2.086 & 43.96\% & 45{,}179 \\
验证 & 绿电比例 & 0.0204 & 12.54\% & 50{,}688 \\
\bottomrule
\end{tabular}
\end{table}

\section{边车怎么接到省平台}
部署形态是策略边车，不是替换省级调度核心。计算放在只读副本旁边：读指标、出建议、写凭证。执行仍走平台现有队列。本仓库的 \path{server.py} 是这个边车的本地、无状态版本，注释写明 Optional stateless sidecar。

本机已经能启动的部分：用 CPU 加载 \path{policy.onnx} 与 \path{forecast.pt}，本次校验在 CPUExecutionProvider 上完成，logit 最大误差 3.81e-06；\path{POST /api/v1/simulate} 在独立内存日历上从种子重放，默认种子 313，小时数 1--168；\path{POST /api/v1/fault/inject} 注入离线、容量损失或 PUE 抬升；\path{POST /api/v1/copilot/command} 解析中文指令并重算；\path{POST /api/v1/ledger} 与 \path{/api/v1/ledger/csv} 导出凭证。

三阶段的进入条件仍是计划，尚未开始。影子阶段要求平台提供至少 14 天只读时序，边车每天出一份与现行策略并列的建议，不对生产队列写指令。人工确认阶段的建议单包含任务标识、目的节点、预计时延圈、演示电价下的电费和凭证哈希；值班员否决时保留否决原因。闭环阶段先在书面变更里限定任务类型和回退开关。生产检查点复制的耗时要单独计量，不能沿用仿真 \path{elapsed_s}。

仓库没有省平台或 13 个地市子门户的接口文档。上述进入条件是对接方案，不是已完成的联调记录。

\section{建议映射、一张真实凭证和上线前要改的默认值}
表~\ref{tab:map} 是边车希望读取的量。右列只说明仓库里用什么字段模拟它。省平台的真实字段名待联调时替换，不能把仿真字段名当成接口规范。

\begin{table}[ht]
\centering
\caption{对接状态。未联调的行保持待联调或待核实。}
\label{tab:map}
\small
\begin{tabularx}{\textwidth}{lXl}
\toprule
边车要读的量 & 仿真里对应的来源 & 对接状态 \\
\midrule
任务类型、到达、时延上限、截止期 & Trace 的五类任务 & 待联调 \\
节点健康、PUE、队列 & \path{world.health} 与功率模型 & 待联调 \\
分时电价 & \path{weather()} 演示曲线；公告价只在小时环境 & 待核实到户电价 \\
绿电可用功率 & 场景参数，不是机房电表 & 待联调 \\
跨域带宽与时延 & 设计矩阵 1/5/10 ms，带宽为假设 & 待核实 \\
回写调度指令 & 无。API 不发生产指令 & 阶段 1 保持不回写 \\
\bottomrule
\end{tabularx}
\end{table}

本次账本样例取自这次 \path{run_benchmark} 写出的 \path{ledger_sample.json}，不是冻结前端包。常规工作日、SuCompute、任务 1：HPC，来源南京，槽 0--14，节点盐城 / DCU，动作概率 0.975，绿电份额 98.07\%，PUE 1.157。有效 1.984 EF$\cdot$h，分摊电量 27{,}600 kWh，演示电价下 10{,}212 元，碳排 293.0 kg，相对全购电少 14{,}887 kg，SLA 满足。SHA-256 \path{e699c16ae9430f01cab21b2efbe628acc9a1f5f99edbd2189af096163dafcc78}。

这张凭证的电量和电费随 10{,}000 kW/EF 放大，不能当成账单或碳资产。\path{server.py} 当前把 CORS 设为允许任意来源，只适合本机演示；接到平台前必须改成指定源，并把 LLM 密钥留在环境变量里，不写入仓库。

\section{验收时对照什么}
\begin{itemize}
\item 权重：\path{verify_weights.py} 通过，五份 SHA-256 与 \path{manifest.json} 一致。
\item 周指标：复跑 \path{run_benchmark.py} 后，四场景等权表与当时的 \path{results_report.md} 一致；墙钟允许变化，并写明主机。
\item 哈希：逐小时 CSV 的 SHA-256 与当次 summary 的 \path{detail_sha256} 一致。不要要求它等于冻结包的 \path{07625d6e37ac223cf1d6b936b6cc3b972d0b16c0ceb32d21cefc413e09437f7d}，除非在同一数值路径上复现到逐小时电量。
\item 账本：抽一条凭证，用 \path{backend.report.receipt_hash} 重算。展示字段 \path{selection_reason} 不进入哈希。
\item 影子试点：连续只读运行，建议单可导出 CSV，生产队列无写入。这项尚未开始。
\end{itemize}

本次复现已经完成前三项里的权重校验和周评测。逐小时哈希等于本次 summary，不等于冻结包。第四项的哈希算法在代码和 \path{API_CONTRACT.md} 里，计划书引用的那条样例来自本次账本文件。第五项没有环境，不能勾。

PPT、录屏和网盘链接不在这份佐证里。优化幅度若没有出现在本次 \path{results_report.md}，就不能写成已验证。

\begin{thebibliography}{9}
\bibitem{report} \path{deliverables/repro/results_report.md} 与 \path{results_summary.json}。命令 \path{python benchmarks/run_benchmark.py --seed 313 --workers 2}。
\bibitem{base} \path{deliverables/repro/baseline_seed313.md}。命令 \path{python scripts/run_baselines.py --seed 313}。
\bibitem{weights} \path{weights/manifest.json}、\path{forecast_metrics.json}、\path{verify_weights.py}。
\bibitem{model} \path{MODEL.md} 与 \path{API_CONTRACT.md}。
\end{thebibliography}
\end{document}
